\section{Controller (ACC, CACC, PLOEG)}
\label{sec:sota_controller}

Il controllo longitudinale costituisce il cuore del platooning: l’Adaptive Cruise Control (ACC) si basa prevalentemente su sensori locali, mentre il Cooperative ACC (CACC) sfrutta informazioni condivise tra veicoli (ad es.\ dal leader o dal predecessore) per migliorare le prestazioni e ridurre le distanze di sicurezza mantenendo stabilità della colonna \cite{ploeg2011design,oncu2014cooperative}. Un obiettivo chiave è la \emph{string stability}, ossia la capacità di non amplificare perturbazioni lungo il platoon, proprietà che dipende in modo critico da politiche di spaziamento e dalla qualità della comunicazione \cite{besselink2017string}. In condizioni di rete degradate, la letteratura evidenzia la necessità di meccanismi di \emph{graceful degradation} e di transizione tra controllori (da CACC ad ACC fino alla guida autonoma non cooperativa), per preservare vincoli di sicurezza anche in presenza di perdite o ritardi \cite{ploeg2015graceful,segata2023multi-technology}. Questa sezione riassume i controllori più rilevanti per il seguito, con particolare attenzione ai presupposti comunicativi e alle implicazioni delle disconnessioni temporanee.
